\section{Conclusions}

The long-only minimum variance portfolio is conventionally solved by forming an $n\times n$
covariance matrix and passing it to a general-purpose solver. This paper shows that the
covariance matrix is unnecessary: the problem reduces to a handful of symmetric positive
definite linear solves (one per balance constraint, plus one for a return tilt, with
the multipliers eliminated through a $p\times p$ Schur complement) and non-negativity
is enforced by an active-set loop. The method works entirely with
the $T\times n$ returns matrix; the covariance matrix is never assembled during
the solve.

The empirical results on 494 S\&P~500 constituents over five years confirm the practical
value of this approach. Against the same solvers invoked through their native APIs
(the fair comparison, free of any modelling-layer overhead), CG is $5.7\times$ faster than
Clarabel without shrinkage and $22\times$ faster under heavy shrinkage ($\alpha=0.5$); the
larger $156\times$ and $661\times$ figures, measured against CVXPY-wrapped Clarabel, additionally
absorb the problem-construction overhead of a general-purpose modelling layer and should be
read as end-to-end rather than algorithmic gains. At large $n$, CG's runtime
grows substantially subquadratically: at $n=3000$ (factor-model synthetic data, $T=1250$ fixed)
CG takes $0.081$\,s, a gap over interior-point methods that widens as $n$ grows
(Figure~\ref{fig:scaling}). Ledoit-Wolf shrinkage benefits all solvers by compressing the
eigenvalue spectrum: it reduces CG iterations from 684 to~82 and OSQP iterations from 225 to~100.
The outer active-set loop also contracts from 18 to~6 steps, because a lower condition number
moves the unconstrained portfolio toward equal weights, so fewer assets receive negative weights
in the primal step. Notably, OSQP called through its direct API
(bypassing CVXPY overhead) is still 6--10$\times$ slower than CG (ranging from $6\times$ without shrinkage to $10\times$ at $\alpha=0.5$):
even a well-engineered general QP solver cannot match algorithms that exploit the
positive-definite structure and operate without assembling the covariance matrix.

Heavy LW shrinkage compresses $\kappa$ as a free by-product of out-of-sample stability
tuning, a connection not previously made explicit in the portfolio optimisation literature.
The statistically optimal $\alpha \approx 0.017$ barely helps conditioning, because the
Frobenius loss is dominated by large, well-estimated eigenvalues that are far from zero;
the near-zero noise eigenvalues, the sole cause of $\kappa \approx 109{,}000$, contribute
negligibly to the loss and receive negligible correction. Practitioners who apply
$\alpha \approx 0.3$--$0.5$ for out-of-sample stability obtain the preconditioning benefit
as a free by-product. RMT eigenvalue cleaning~\cite{schmelzer2026rmt} resolves the
tension more decisively by clipping noise eigenvalues to their theoretically predicted
bulk mean, reducing $\kappa$ to~146 on S\&P~500 data; the Woodbury direct solver for the
RMT target is developed in the companion paper.

The projected-gradient baseline of the companion paper~\cite{schmelzer2026nncg} is
$19\times$ slower than CG at $n=3000$,
as its $O(\kappa)$ iteration count grows with the condition number while CG's $O(\sqrt\kappa)$
count does not. The method extends naturally to the
mean-variance problem with a return tilt, which requires one additional SPD solve with
the balance multipliers still recovered analytically, as outlined in Section~\ref{sec:pd}.
The frontier experiment in Section~\ref{sec:results} demonstrates this concretely:
a 50-point long-only efficient frontier for $n=500$ assets is completed
$1{,}455\times$ faster than CVXPY (warm-started CG with LW shrinkage, Table~\ref{tab:frontier}).
